\documentclass[11pt]{article}

\usepackage[margin=1in]{geometry}
\usepackage{amsmath}
\usepackage{amssymb}
\usepackage{amsthm}
\usepackage{booktabs}
\usepackage{array}
\usepackage{enumitem}
\usepackage{xcolor}
\usepackage{hyperref}
\hypersetup{colorlinks=true, linkcolor=blue, urlcolor=blue, citecolor=blue}

\newcommand{\code}[1]{\texttt{#1}}
\newtheorem{definition}{Definition}
\newtheorem{proposition}{Proposition}
\newtheorem{remark}{Remark}

\newcolumntype{L}[1]{>{\raggedright\arraybackslash}p{#1}}

\tolerance=800
\emergencystretch=3em

\title{\textbf{Attested Inference on Operator Hardware}\\[2mm]
\large What confidential computing establishes for a third-party AI operator,\\
what the settlement commitment adds, and what remains unsolved}

\author{Hanzo AI Research\thanks{Corresponding author: research@hanzo.ai} \\
    \textit{Hanzo Industries Inc (Techstars '17)} \\
    Los Angeles, California \\
    \texttt{https://hanzo.ai}}

\date{Revision~1}

\begin{document}
\maketitle

\begin{abstract}
Hanzo intends to buy inference capacity from operators it does not own, and to pay them for it
automatically. Two properties have to hold before that is safe: a user's prompt must stay
confidential on a machine its owner controls, and the operator must not be able to bill for a model
it did not run. Confidential computing is frequently offered as the answer to both. It is the answer
to neither in full, and the two failures are of different kinds.

A hardware attestation from AMD SEV-SNP, Intel TDX, or an NVIDIA GPU in confidential mode
establishes that an environment whose initial memory hashes to a stated measurement was launched on
a genuine device of a stated family, at a stated firmware version, and that a 64-byte field
originated inside that environment. It establishes nothing about the behaviour of the software so
measured; a measurement is a hash of bytes, not a proof about what those bytes do. A model
commitment --- a digest over the weights actually loaded --- is similarly a membership statement: it
binds a name to bytes, and does not establish that those bytes were the ones multiplied. The claim
``the advertised model produced this output'' is therefore not proved by either mechanism
separately, nor by both together. It is \emph{inherited} from a chain of assumptions about hardware,
firmware, and the correctness of measured software, and it fails silently if any link fails.

This document specifies the flow anyway, because inheriting a property under named assumptions is
better than asserting it without them, and because the mechanisms that do not rest on the operator's
hardware --- an exact-arithmetic product check, and a sampled audit against the reference weights
--- compose with it rather than replacing it. It states what each layer proves; classifies every
resulting guarantee as cryptographic, hardware-trust, economic, or open; and reports the observed
state of the implementation in the Hanzo and Lux trees, including three respects in which the
attestation verifier does not yet supply what the on-chain evidence backends consuming it assume.
\end{abstract}

\tableofcontents

\section{Scope}

The setting is an operator who owns a machine, wants to earn from it, and is not Hanzo. The
operator installs a Hanzo image, the platform routes inference to it, and the platform pays for the
work. HIP-0096 \cite{hip0096} specifies the reward economics for that arrangement and HIP-0095
\cite{hip0095} the quality-of-service challenges around it; HIP-0901 \cite{hip0901} specifies Proof
of AI, the settlement proof system, and assigns confidential compute to the tier where exact
arithmetic is unavailable. The machine an operator connects, and the socket it holds, are HIP-1325's
subject \cite{hip1325}. What
none of them specifies is the confidential-compute tier itself: which evidence is collected, which
fields a verifier checks, what the resulting statement means, and where it stops being a proof.

Two properties are at stake, and they must not be conflated.

\begin{description}[leftmargin=1.2em,itemsep=3pt]
\item[Confidentiality.] The user's prompt, the returned completion, and (where the weights are not
public) the model itself are not readable by the operator who owns the machine they run on.
\item[Execution integrity.] The output the operator returns and bills for is the output the
advertised model produces on the given input under the advertised decoding regime.
\end{description}

Confidential computing addresses the first directly and the second only by implication. The
implication is the subject of most of this document.

Two things are out of scope. Training and fine-tuning on operator hardware raise a different set of
questions (gradient poisoning, data extraction from updates) and are covered by HIP-0901's
contribution proof. Availability --- an operator who simply stops answering --- is a routing and
reputation matter, not a cryptographic one, and is handled by HIP-0095.

\section{Vocabulary}

Every term below is used in exactly one sense for the remainder of the document.

\begin{definition}[Measurement]
A cryptographic digest over the initial state of an isolated execution environment: the memory pages
loaded before the environment began executing, together with the metadata that describes how they
were loaded. Written $M$. On SEV-SNP it is the 48-byte launch digest; on TDX it is \code{MRTD}
together with the runtime measurement registers \code{RTMR0..3}, which record what the environment
loaded after launch.
\end{definition}

\begin{definition}[Attestation]
A statement signed by a hardware root of trust asserting that an environment with measurement $M$,
running under configuration $c$ (security version, policy bits, firmware version), was instantiated
on a genuine device of a stated family, and that a caller-supplied field $\mathrm{rd}$ (64 bytes on
both SEV-SNP \cite{snpabi} and TDX \cite{tdx}) was provided by software executing inside that
environment.
\end{definition}

\begin{definition}[Trusted computing base]
The set of components whose correct behaviour is necessary for a stated property to hold. Written
TCB. A component is in the TCB whether or not anyone examined it.
\end{definition}

\begin{definition}[Model commitment]
A digest $D_w$ over the weight tensors as loaded, including a per-tensor discriminant for the
quantization scheme, such that re-quantizing the same model yields a different digest. HIP-0901
calls this \code{modelSpecHash} and computes it as a keccak Merkle root over the quantized weight
bytes.
\end{definition}

\begin{definition}[Receipt]
A signed statement, emitted by the software inside the environment, binding the model commitment,
the decoding regime, the input, and the output of one request. Written $R$.
\end{definition}

\begin{definition}[Execution integrity]
The property that a returned output equals the output of the advertised model on the given input
under the advertised decoding regime. This is a statement about a computation, not about an
environment.
\end{definition}

\begin{remark}
Attestation, model commitment, and execution integrity are three distinct properties. Attestation is
about an environment. A model commitment is about bytes. Execution integrity is about a computation.
No amount of the first two entails the third without additional assumptions, which
Section~\ref{sec:notproved} states.
\end{remark}

\section{What the tree contains}
\label{sec:observed}

This section reports the implementation as observed, so that later sections can distinguish what is
running from what is proposed. Line counts and test counts are from direct inspection.

\subsection{Off-chain: privacy tiers and key brokerage}

Two Rust crates model the confidential-compute surface.
\code{net/hanzo-pqc/src/privacy\_tiers.rs} defines a five-level ladder --- open data, at-rest
encryption, CPU TEE, CPU TEE with GPU confidential computing, CPU TEE with GPU TEE-I/O --- together
with a capability structure that derives the maximum tier a host supports from its vendor features.
\code{net/hanzo-pqc/src/attestation.rs} declares a verifier interface with two operations (verify a
CPU quote against an expected measurement; verify a GPU token) and a field-by-field decoder for the
1184-byte SEV-SNP report structure.

The verifiers that implement that interface are mocks. \code{MockAttestationVerifier} in
\code{hanzo-pqc} returns success unconditionally for any quote and any expected measurement, and
returns a fixed H100 result for any GPU token; the production steps are enumerated in comments. The
sibling crate \code{rust-sdk/crates/hanzo-kbs} carries the key-management and key-brokerage split
--- lifecycle in one module, attestation-gated release in the other --- and its
\code{MockAttestationVerifier} accepts any non-empty report as an SEV-SNP attestation at tier
\emph{CpuTee}. Neither crate contains a certificate-chain walk, a measurement allowlist, a policy-bit
check, or a TCB floor.

The two crates also define two overlapping enumerations for the same concept: \code{PrivacyTier} with
five levels ordered by protection, and \code{CCTier} with four levels ordered inversely and named for
mechanism rather than for what is protected. One of them should survive.

\subsection{Node-local: boot measurement and host-asserted evidence}
\label{sec:hanzoattest}

\code{hanzo-libs/hanzo-attest} in \code{hanzoai/node} supplies the layer below this document. It
defines a boot measurement --- a digest, under a canonical encoding, over the inputs that determine a
microVM: the monitor binary, its operating-system asset version, the kernel image, the root
filesystem, the machine shape, and the guest argv --- and an evidence type that is a signed statement
about one such measurement. The signing preimage carries a domain tag, a byte identifying what roots
the statement, the measurement, and the second it was made, so a statement cannot be replayed later
under a fresher date and a signature made under one root cannot be presented as another.

Its enumeration of roots is the boundary this document begins at. \code{Host} --- the node's own
identity key --- is implemented, and establishes that the statement came from that node and was not
altered afterwards; it does not establish that the node fed those files to the monitor, because
nothing outside the node observed the boot. \code{Snp} and \code{Tdx} are declared and refused by
name rather than accepted silently: there is no attester and no verifier for either. The companion
paper \cite{hanzodattest} specifies that measurement and its trust base in full. Everything below
assumes a hardware root exists and asks what follows from it.

\subsection{On-chain: evidence backends}

\code{luxfi/standard} at \code{contracts/ai/compute/} carries a verifier that dispatches to evidence
backends by proof type, and a profile contract mapping a task tier to the proof type it requires.
Two of the backends are confidential-compute backends. \code{CCEvidence} (proof type 1) accepts a
\code{reportData} if a governance-accepted attestation root has vouched for it. \code{TEEEvidence}
(proof type 4) requires a threshold number of distinct governance-approved enclave measurements to have each vouched
for the same \code{reportData}, deduplicating repeated entries so one enclave cannot fill the quorum
alone.

Both are fail-closed in the deployed configuration, and deliberately so: the voucher set is empty
and \code{TEEEvidence}'s threshold defaults to zero, which the contract treats as ``never attest''.
Nothing settles at these tiers today. \code{TEEEvidence}'s own header states the trust class
plainly, describing the tier as an attested-execution tier rather than a cryptographic-soundness
tier, and identifying the quote signature as the one classical elliptic-curve assumption in a system
that is otherwise post-quantum. That is the correct classification and this document adopts it.

\subsection{On-chain: what the attestation precompile actually verifies}
\label{sec:precompile}

The heavy verification is delegated off the EVM. \code{luxfi/precompile} carries an attestation
precompile at \code{0x0301}, a stateless adapter over one verifier, \code{ai.VerifyTEE}. That
verifier does two things: it parses a concatenated DER certificate chain and requires the leaf to
chain to an embedded vendor root, valid at the timestamp carried in the report; and it checks that
the leaf's certified key signed the report. Verification is deterministic --- the chain is validated
at the report's own embedded timestamp rather than at wall-clock time, so every validator reaches
the same verdict --- and it fails closed on an unparseable chain, an untrusted root, or a signature
that does not bind the report.

Three properties of that verifier bear directly on this document, and all three are observations
about the current code rather than criticisms of its design intent.

First, the object it verifies is not an SEV-SNP report or a TDX quote. It is a Lux-defined 48-byte
header: a 32-byte device identifier, an 8-byte big-endian timestamp, and an 8-byte nonce, followed
by the certificate chain. There is no measurement field, no policy word, no reported TCB version,
and no 64-byte report-data field. A successful verification therefore establishes that \emph{a key
certified under an embedded vendor root signed these 48 bytes}, and no more.

Second, \code{TEEEvidence}'s documented contract with its off-chain voucher requires the voucher to
check that ``the quote's report-data field equals \code{reportData}'' before recording a vouch. The
payload the current verifier signs over has no such field. The binding between an attestation and
the computation it is supposed to witness is, at present, not carried by the verified bytes.

Third, exactly one vendor root is embedded: the Intel SGX Provisioning Certification Root CA. AMD
SEV-SNP chains, which terminate at an AMD root, and NVIDIA device-identity chains both fail closed
today. The code says so, and the pool is explicitly never the host's system store.

The freshness question follows from the first point. The header carries a nonce, but nothing in the
verifier compares it against a challenge, and the chain is checked at a timestamp the report itself
asserts. A report is therefore replayable indefinitely by anyone holding a copy unless a caller
enforces the nonce out of band. Section~\ref{sec:evidence} makes that enforcement explicit.

\subsection{The proof system this composes with}

HIP-0901 specifies Proof of AI: a Freivalds \cite{freivalds1977} matrix-product check over the exact integer accumulator
of the engine's int8 path, in the Mersenne prime field $p = 2^{61}-1$, with soundness error at most
$p^{-k}$ for $k$ challenge vectors. The Go verifier at \code{luxfi/crypto/poi} is present with 22
test functions across its Freivalds, transcript, wire, scale, and adversarial suites. The Rust
prover core at \code{hanzo-engine/src/poi.rs} is present with 7 \code{\#[test]} functions in that
file and 38 across the proof-of-inference surface as a whole.\footnote{HIP-0901 states 14
\code{\#[test]} in \code{poi.rs}. The tree shows 7 there and 38 across \code{poi.rs},
\code{poi\_transcript.rs}, \code{poi\_graph.rs}, \code{poi\_forward.rs}, and
\code{tests/poi\_adversarial.rs}. The HIP's figure should be reconciled against the tree; the
substance of its claim --- that the core is implemented and adversarially tested --- holds.} This
matters here because Freivalds is the one path in the whole system that yields a cryptographic
statement about execution, and its applicability is restricted to exact integer arithmetic
(Section~\ref{sec:freivalds}).

\subsection{Hardware available for this work}

Of the two hosts on the development bench, one was reachable at the time of writing. It reports as
\code{spark}, aarch64, Ubuntu 24.04.4, kernel 6.17.0-1031-nvidia, with an NVIDIA GB10 and driver
580.173.02. It exposes no CPU TEE: \code{/proc/cpuinfo} carries no SEV, SEV-ES, SEV-SNP or TDX flag
(it is an Arm part), \code{/dev/sev} is absent, the kernel log mentions no realm or CCA support, and
the installed \code{nvidia-smi} has no \code{conf-compute} subcommand and reports no
confidential-computing section under \code{-q}. The second host, an amd64 machine, did not answer.

The flow specified below is therefore \emph{specified and not exercised}. No measurement in this
document has been taken on confidential hardware, and no performance claim is made about
confidential mode.

\section{The confidential path}
\label{sec:flow}

The path has four phases. They are named rather than numbered because they are not a sequence a
reader follows once; enrollment happens once per operator, launch once per boot, release once per
key lifetime, and a session once per user connection.

\subsection{Image and measurement}

Hanzo publishes an image, its digest, and the launch measurement $M$ that image produces under a
stated launch configuration. The measurement is a function of the image bytes, the guest firmware,
the number of vCPUs, and the launch policy; changing any of them changes $M$.

A measurement is meaningful only if a reader can map it back to source. This is the requirement most
easily skipped and the one that makes the difference between an attestation that means something and
an attestation that means ``the same unknown thing as last time''. Three things are needed together:
the image build must be reproducible, so that an independent party rebuilding from the named source
revision obtains the same bytes; the mapping from image digest to measurement must be computable
offline, so a verifier can derive the expected $M$ rather than being told it; and the set of
$(\text{revision}, \text{digest}, M)$ triples must be published in an append-only log with the
lone-node promotion rule of RFC 6962 \cite{rfc6962}, so that a measurement quietly added to an
allowlist is visible and a subtree cannot be reinterpreted as the whole log.

Without those, \code{approveMeasurement} in \code{TEEEvidence} is an assertion by whoever holds the
admin key that a 32-byte value is trustworthy, and the quorum of distinct measurements it
enforces is a quorum over values nobody outside can interpret.

\subsection{Launch and attestation}

The operator boots the measured image in a confidential VM. Inside, before serving anything, the
guest generates two ephemeral key pairs in guest memory and persists neither:

\begin{itemize}[leftmargin=1.4em,itemsep=2pt]
\item $(pk_e, sk_e)$, an X-Wing key pair --- the hybrid of X25519 and ML-KEM-768 specified under
that name ---
used to receive sealed prompts and sealed key material, under HPKE \cite{hpke}. HIP-0901 already
specifies X-Wing for prompt sealing at \code{crypto/promptseal}; this reuses it rather than
introducing a second envelope.
\item $(pk_r, sk_r)$, a receipt signing key.
\end{itemize}

The guest then requests a hardware report with the report-data field bound to both keys, the session
identifier, and a verifier-supplied nonce:
\[
\mathrm{rd} \;=\; \mathsf{H}\big(\text{``hanzo/attest/v1''} \,\|\, pk_e \,\|\, pk_r \,\|\, \mathrm{sid} \,\|\, n_v\big),
\]
where $\mathsf{H}$ is keccak-256, matching HIP-0901's hash so the system has one. The evidence
bundle sent to the verifier is the report, its certificate chain, $pk_e$, $pk_r$, $\mathrm{sid}$,
$n_v$, and --- where the tier requires it --- one GPU attestation report per device, each generated
with a GPU nonce derived as $\mathsf{H}(\text{``hanzo/attest/gpu/v1''} \| \mathrm{rd} \| \text{device index})$.
Deriving the GPU nonce from the CPU report data is what ties the two attestations into one
statement; without it, an operator can present a CPU report from one machine and a GPU report from
another.

The two-level key structure is deliberate. Hardware report generation is comparatively expensive and
rate-limited, so it happens once per boot; $sk_r$ then signs per-request receipts at negligible
cost. The consequence, stated plainly, is that everything after the attestation rests on $sk_r$
remaining inside the environment. Section~\ref{sec:residual} treats its extraction as a live
possibility rather than an edge case.

\subsection{Key release}

Where the model weights are not public, they must reach the operator without the operator being able
to read them. The KMS at \code{kms.hanzo.ai} releases a wrapped data-encryption key to $pk_e$, sealed
with X-Wing HPKE, only when every condition in Section~\ref{sec:evidence} holds. The release is
scoped to the model set the operator is licensed to serve, carries a short lifetime after which the
environment must re-attest, and is recorded in the KMS audit trail as an access by a measured
environment rather than by a human.

The enclave decrypts the weights into guest memory and never writes them back in plaintext. This is
the strongest available arrangement and it is not a guarantee. An adversary who breaks the
environment's isolation obtains the weights, and Section~\ref{sec:trust} explains why that adversary
is unusually well positioned here: he owns the machine. The architectural consequence is a policy
one --- weights whose disclosure would be catastrophic do not belong on third-party operators at any
tier, and the tier ladder should be read as a statement about cost to the attacker, not about
impossibility.

\subsection{Session}

The user's client obtains $pk_e$ together with the evidence bundle and verifies the evidence itself.
This is the point at which the platform stops being trusted: a client that accepts $pk_e$ on the
platform's word has delegated the entire question to Hanzo, which is a defensible product decision
and a different security property from the one this document specifies. The verification is
reproducible offline from published measurements and vendor roots, so a client can perform it.

The client seals the prompt to $pk_e$ using promptseal, with the additional authenticated data
binding the request identifier so a sealed prompt cannot be replayed under a different request. The
environment opens it, runs the model, seals the completion back to the client, and emits the receipt
of Section~\ref{sec:receipt}. The operator sees ciphertext in, ciphertext out, and metadata
throughout.

\subsection{The device boundary, which is where this usually breaks}

A CPU TEE protects memory the CPU owns. A discrete GPU is not inside it. A device that is not itself
a trusted execution environment cannot perform DMA into encrypted guest memory, so transfers between
a confidential VM and a passthrough GPU pass through shared pages, and the tensors crossing that
boundary --- weights on load, activations and KV cache in flight --- are readable by the host.

Stated directly: \emph{a CPU TEE wrapped around a passthrough GPU protects the control plane and
not the computation.} At tier 2 in the ladder of
Section~\ref{sec:observed}, prompts are confidential against a passive host only while they are on
the CPU. Confidentiality of the inference itself begins at tier 3, where the GPU runs in
confidential mode: the device refuses external access to its protected memory regions, and DMA
between the CPU environment and the device travels over an encrypted, integrity-protected channel
established through the device's own identity chain. Tier 4 (TEE-I/O) removes the bounce buffers by
admitting the device into the trust domain directly.

Two costs follow. The trust base grows to include the GPU vendor, its firmware, and its attestation
service (Section~\ref{sec:trust}). And confidential-mode overhead falls almost entirely on the
host-device transfer path, so it depends on the ratio of transfer to compute in a given workload; we
have taken no measurement of our own and make no claim about the magnitude.

\section{The evidence a verifier must check}
\label{sec:evidence}

This section is a checklist, and it is the part of the design where omissions are silent. Every item
below has been the subject of a published break or a documented misuse in some deployment; a
verifier that skips one does not fail loudly, it accepts.

\paragraph{AMD SEV-SNP.}
Validate the certificate chain from the VCEK, through the ASK, to the ARK, for the chip identifier
and the reported TCB version in the report itself --- the VCEK is derived from both, so presenting a
report whose \code{reported\_tcb} does not match the key that signed it must fail. Check
\code{MEASUREMENT} against the published allowlist. Check the policy word: \code{DEBUG} must be
clear, or the host may attach a debugger to the guest and the entire construction is void;
\code{MIGRATE\_MA} must be clear unless a migration agent is intended; \code{SMT} must match
deployment policy; the minimum ABI version must be met. Require \code{reported\_tcb} and
\code{GUEST\_SVN} to meet a floor, so an operator cannot roll firmware back to a version with a
known break. Check \code{REPORT\_DATA} equals the expected binding. Where an ID block is used, check
\code{ID\_KEY\_DIGEST} and \code{AUTHOR\_KEY\_DIGEST}. Distinguish VCEK-signed from VLEK-signed
reports and accept only what policy allows. Check revocation.

\paragraph{Intel TDX.}
Verify the quote against Intel's provisioning collateral: the quoting enclave identity, the TCB
info, and the revocation lists. Check \code{MRTD} and the runtime measurement registers
\code{RTMR0..3} against the allowlist. Check the TD attributes: \code{DEBUG} clear, \code{XFAM}
within policy. Treat the TCB status strictly --- a status of \emph{configuration needed} or
\emph{software hardening needed} is not \emph{up to date}, and accepting it without a recorded
policy decision silently admits a platform with a known issue. Check the report data binding.

\paragraph{NVIDIA GPU in confidential mode.}
Verify the device attestation report against NVIDIA's device-identity chain \cite{nvcc}, and the reported
firmware and VBIOS measurements against the published reference integrity manifests for that
device and driver. Require confidential mode to be genuinely on rather than in a development or
devtools setting. Check the GPU nonce equals the value derived from the CPU report data. Where a
partitioned device is used, require the partition configuration to appear in the attested state
rather than being reported by the host.

\paragraph{Freshness, in all cases.}
A hardware report carries no trustworthy notion of ``now''. SEV-SNP reports carry no timestamp at
all; the Lux payload of Section~\ref{sec:precompile} carries one the report itself asserts. Freshness
comes only from the verifier's own nonce inside the report data, and from the fact that the keys
bound there are ephemeral and die with the environment. A verifier that does not generate $n_v$
itself, or does not check it, has accepted a report that may be arbitrarily old --- and since a
report is public evidence, an old one is not hard to obtain.

\paragraph{Uniqueness.}
Record the chip identifier and report identifier. One physical machine presenting itself as several
independent operators is the Sybil case HIP-0096's reward economics assumes away; the attestation is
what makes that assumption checkable, and only if the identifiers are recorded and compared.

\section{The settlement receipt}
\label{sec:receipt}

The receipt is what settlement pays against. Its purpose is narrow: to make model substitution a
detectable event rather than an invisible one. It reuses HIP-0901's transcript fields so that the
system has one commitment format, and adds what that binding does not cover: a tokenizer digest, the
sampling seed, and digests of the input and output in the form the user sees them.

For one request, the environment signs
\[
R \;=\; \big(D_w,\; Q,\; D_t,\; \kappa,\; s,\; x_{\text{raw}},\; x_{\text{tok}},\; y_{\text{tok}},\; y_{\text{raw}},\; n_{\text{in}},\; n_{\text{out}},\; \mathrm{sid},\; i\big),
\]
\[
\rho \;=\; \mathsf{H}(\text{``hanzo/receipt/v1''} \,\|\, R), \qquad \sigma \;=\; \mathrm{Sign}_{sk_r}(\rho),
\]
with fields as follows.

\begin{center}
\small
\begin{tabular}{@{}lL{0.72\textwidth}@{}}
\toprule
Field & Meaning \\
\midrule
$D_w$ & Model commitment: keccak Merkle root over the quantized weight bytes as loaded, with the
quantization discriminant per leaf. HIP-0901's \code{modelSpecHash}. \\
$Q$ & Quantization specification. \\
$D_t$ & Tokenizer digest. \emph{New.} \\
$\kappa$ & Engine and kernel version, including the reduction-order and batch-invariance settings in
force. \\
$s$ & Sampler specification: temperature, top-$p$, top-$k$, penalties, \emph{and the seed}. \emph{Seed
is new.} \\
$x_{\text{raw}}$ & Digest of the input as the user sent it, before tokenization. \emph{New.} \\
$x_{\text{tok}}$ & Digest of the canonical input token identifiers. HIP-0901's \code{inputCommitment}. \\
$y_{\text{tok}}$, $y_{\text{raw}}$ & Digests of the output token identifiers and of the detokenized bytes. \\
$n_{\text{in}}$, $n_{\text{out}}$ & Token counts, which is what billing multiplies. \\
$\mathrm{sid}$, $i$ & Session identifier and a counter monotone within the session. \\
\bottomrule
\end{tabular}
\end{center}

\paragraph{Why the tokenizer digest.}
HIP-0901 folds the decoding regime into the quantization specification and the kernel version, and
commits the input as canonical token identifiers. A tokenizer swap changes the map from user text to
token identifiers and therefore changes the output, while leaving the model commitment, the
quantization specification, and the kernel version untouched. The committed token identifiers do
change --- but they are the identifiers \emph{the environment chose}, so a user checking the
commitment against their own text needs to know which tokenizer to use. Binding $D_t$ and
$x_{\text{raw}}$ closes that: anyone holding the advertised tokenizer can recompute $x_{\text{tok}}$
from $x_{\text{raw}}$ and reject a mismatch.

\paragraph{Why the sampling seed.}
Committing the seed converts a stochastic output into a deterministic function of weights, input,
and seed. That is what allows an auditor holding the weights to replay a request exactly rather than
testing it statistically (Section~\ref{sec:audit}), and it costs one field. The seed must be
derived, not chosen: $s.\mathrm{seed} = \mathsf{H}(\mathrm{sid} \,\|\, i \,\|\, n_v \,\|\, n_c)$,
where $n_c$ is a nonce the client places inside the sealed request. Deriving it this way ties the
seed to material fixed before generation and gives the client a term of its own, so an environment
can neither search seeds for a favourable output nor reuse a session's randomness across requests.

\paragraph{Settlement rule.}
Payment is for what the receipt establishes, not for what the operator claims. A request whose
receipt is absent, unsigned by an attested $pk_r$, or bound to a model commitment outside the
operator's licensed set settles at zero. A receipt naming a smaller model than the one advertised
settles at that smaller model's rate. The chain-side path already exists: the receipt hash is the
\code{reportData} that \code{CCEvidence} or \code{TEEEvidence} judges, and HIP-0901's
\code{AIComputeRegistry} keys its no-double-mint ledger on the computation commitment, so one unit
of work settles once network-wide.

\paragraph{The gap this leaves in the current tree.}
Section~\ref{sec:precompile} observed that the payload the attestation verifier signs over has no
report-data field. The receipt hash $\rho$ is exactly what belongs there. Until the verified payload
carries it --- either by parsing a real SEV-SNP or TDX report and reading its 64-byte field, or by
extending the Lux header --- the vouch that \code{TEEEvidence} records cannot mean what
\code{TEEEvidence}'s documentation says it means. The contracts are fail-closed, so nothing is
currently mis-settling; the work is unfinished, not broken.

\section{The trust base}
\label{sec:trust}

\begin{center}
\small
\begin{tabular}{@{}L{0.26\textwidth}L{0.36\textwidth}L{0.30\textwidth}@{}}
\toprule
Component & Trusted for & If it fails \\
\midrule
CPU vendor root and endorsement keys & Authenticity of every report from that vendor's parts &
Arbitrary reports forge arbitrary environments; no downstream check detects it \\
CPU firmware / TDX module at the attested version & The isolation the report asserts &
Confidentiality and integrity of the environment; receipts become operator-controlled \\
GPU vendor identity chain and firmware (tier $\geq 3$) & Device authenticity and confidential mode &
Tensors readable by the host; GPU attestation forgeable \\
The measured image & Implementing the receipt discipline correctly & Receipts are well-formed and
meaningless \\
Reproducible build and measurement log & That $M$ corresponds to reviewable source & Allowlisted
measurements are uninterpretable; the quorum is over opaque values \\
KMS release policy & Releasing weights only to environments meeting the policy & Weights to an
unmeasured host \\
Governance of the measurement set and quorum & Which environments count & An approved malicious
measurement satisfies the quorum trivially \\
Physical integrity of the machine & Everything above it & See below \\
\bottomrule
\end{tabular}
\end{center}

The last row is the one that makes this setting different from a cloud deployment. In the standard
threat model for confidential computing, the adversary is a privileged software attacker: a
malicious hypervisor, a compromised host operating system, a co-tenant. Physical attacks are
acknowledged and largely deferred, on the reasonable ground that a hyperscaler's data centre
restricts physical access.

Here the adversary owns the machine, has unlimited time with it, and is financially motivated by the
system's own reward schedule. Every attack requiring physical access moves from the deferred column
to the live column. Published work on this class includes voltage fault injection against the AMD
secure processor to extract endorsement material \cite{buhren2021}, and memory-aliasing attacks that
defeat SEV-SNP's replay protection using inexpensive hardware to tamper with a memory module's
declared geometry \cite{badram}. Both were addressed by firmware or platform updates for the
affected generations; neither is the last of its kind, because the underlying asymmetry --- the
attacker holds the hardware --- does not change.

The honest summary is that the tier ladder measures \emph{cost to the attacker}, and that for a
sufficiently motivated operator on hardware he owns, that cost is finite and in some published cases
small.

\section{What an attestation does not establish}
\label{sec:notproved}

\begin{proposition}[What a valid report gives you]
\label{prop:report}
Assume: \textup{(H1)} the vendor's endorsement key is uncompromised and the certificate chain
validates; \textup{(H2)} the hardware and firmware at the attested version enforce the isolation
they specify; \textup{(H3)} the verifier performs the checks of Section~\ref{sec:evidence},
including the freshness binding. Then the verifier learns that an environment whose initial state
hashes to $M$, under configuration $c$, was instantiated on a genuine device of the stated family,
and that $\mathrm{rd}$ was supplied from inside it. It learns nothing further.
\end{proposition}

\begin{proposition}[A measurement carries no semantics]
\label{prop:semantics}
$M$ is a digest of initial memory. Two images differing in one instruction have unrelated
measurements; an image and its malicious variant are equally well attested. A verifier who cannot map
$M$ to reviewable source learns only that the environment is the same one as before, not what it
does.
\end{proposition}

\begin{proposition}[Execution integrity is inherited, not proved]
\label{prop:inherit}
Assume additionally: \textup{(H4)} the software with measurement $M$ implements the receipt
discipline of Section~\ref{sec:receipt} correctly --- it loads exactly the weights whose digest it
reports, applies exactly the decoding regime it reports, and reports the output it produced; and
\textup{(H5)} no adversary extracts $sk_r$ or influences the environment's control flow or data.
Then a receipt verifying under $pk_r$ implies that the advertised model produced the reported output
on the reported input.

The conclusion rests on \textup{(H1)--(H5)}. None of them is a cryptographic assumption about the
receipt itself: the signature establishes only that the holder of $sk_r$ signed $\rho$. Drop
\textup{(H4)} and the receipt is a well-formed statement about nothing. Drop \textup{(H5)} and it is
a statement the operator can author at will.
\end{proposition}

\begin{proposition}[A commitment establishes membership, not use]
\label{prop:commitment}
Let $D_w$ be a digest over the advertised weights. That $D_w$ is the advertised model's digest is
publicly checkable by anyone holding the weights. That the bytes under $D_w$ were the operands of
the multiplications producing $y$ is not established by the commitment, and is not established by
the attestation; under Proposition~\ref{prop:inherit} it is inherited from \textup{(H4)} and
\textup{(H5)}.
\end{proposition}

\begin{proposition}[The one cryptographic exception]
\label{prop:freivalds}
Let a transcript commit to the operands of a layer's product and let each committed operand carry a
Merkle inclusion proof to $D_w$. Let a challenge vector $r$ be derived from a beacon fixed after the
transcript is committed. Then the Freivalds check $A(Br) = Cr$ over $\mathbb{F}_p$ fails with
probability at least $1 - p^{-1}$ per vector if $C \neq AB$. This is a statement about the
computation and it does not depend on \textup{(H1)--(H5)}. Its soundness is information-theoretic,
so it holds against an unbounded adversary, conditional on the beacon being unpredictable at the
moment the transcript is committed --- a liveness property of the chain, not a computational
hardness assumption.

What the check establishes is that the committed $C$ is the product of the committed $A$ and $B$;
the inclusion proof to $D_w$ is what makes $A$ the advertised weights rather than arbitrary bytes.
Its scope is bounded by its arithmetic, which must be an exact integer accumulator, and by what is
opened --- it says nothing about layers left unopened, which is why HIP-0901 pairs it with bisection
to the first divergent layer.
\end{proposition}

Proposition~\ref{prop:freivalds} is the only cryptographic statement about execution available in
this system, and it applies to the quantized integer path. For floating-point inference --- which is
where most GPU serving happens --- there is no equivalent, and confidential computing is the
fallback precisely because nothing stronger exists at that numeric tier.

\section{What a malicious operator retains}
\label{sec:residual}

Assume the flow of Section~\ref{sec:flow} is deployed correctly and the checks of
Section~\ref{sec:evidence} are all performed. The operator can still do all of the following.

\paragraph{Refuse, delay, and select.}
Nothing here compels service. An operator can drop requests, serve slowly, serve well only when it
suspects observation, or refuse specific users identifiable by network metadata. Availability and
consistency are HIP-0095's subject; the point here is that confidentiality and integrity mechanisms
do not touch them.

\paragraph{Read the metadata.}
Prompt and completion are sealed; their lengths, timings, arrival rates, the identity of the
connecting client, and the sequence of token-generation intervals are not. Inter-token timing in
particular is a side channel of practical significance for a streaming interface. Padding and
batching reduce it at a cost in latency; the current design does neither, and this document does not
claim metadata privacy.

\paragraph{Attack the environment from below.}
The operator controls the hypervisor, the interrupt controller, the page tables outside the
protected set, and the power supply. The published literature against SEV-SNP and TDX includes
ciphertext side channels arising from deterministic memory encryption \cite{cipherleaks}, precise
single-stepping frameworks that turn a coarse channel into an instruction-level one
\cite{sevstep}, malicious interrupt and exception injection into a protected guest
\cite{heckler,wesee}, and cache-writeback manipulation to revert protected memory \cite{cachewarp}.
Each has a mitigation in current firmware; the class does not close, and the enumeration above is
illustrative rather than complete.

\paragraph{Extract $sk_r$, once.}
This is the highest-value target in the design. A single successful
extraction of the receipt signing key --- by side channel, fault injection, or a bug in the measured
image --- lets the operator author valid receipts for computations it never performed, indefinitely,
with no observable difference from honest operation. Attestation cannot detect it, because the
attestation was genuine when it was made. The mitigations are bounded key lifetime with mandatory
re-attestation, the distinct-measurement quorum that \code{TEEEvidence} already
implements, and the independent audit of Section~\ref{sec:audit} --- which is the only one of the
three that detects the compromise rather than limiting its window.

\paragraph{Roll back, if allowed.}
Firmware and security-version numbers appear in the report. An operator running an older version
with a published break is detectable, and detectable only if a minimum is enforced. Without that
minimum every other check in Section~\ref{sec:evidence} still passes.

\paragraph{Substitute upstream of the environment.}
Everything above concerns the operator's ability to corrupt what happens inside. The operator also
controls what reaches it. A request routed to a cheaper endpoint that never enters the environment
produces no receipt at all --- which the settlement rule of Section~\ref{sec:receipt} handles by
paying zero, provided the platform actually requires a receipt for every billed request and does not
fall back to a claimed usage record when one is missing. Such a fallback is easy to add for
operational reasons and removes the rule's effect entirely, so it should be refused explicitly
rather than left to convention.

\section{Verification that does not rest on the hardware}
\label{sec:audit}

Four mechanisms give evidence about execution without assuming the operator's hardware behaves.
They differ in what they establish and in cost.

\subsection{Exact arithmetic: Freivalds}
\label{sec:freivalds}

Proposition~\ref{prop:freivalds} gives a genuine cryptographic check where the accumulator is exact
integer arithmetic. HIP-0901 specifies it, \code{luxfi/crypto/poi} implements the verifier, and
\code{hanzo-engine} implements the prover core. It requires the served path to be the quantized
integer path; it does not extend to bf16 or fp8 serving, and a tolerance-based variant would be both
unsound (a cheat hides under the tolerance) and false-rejecting (honest reductions differ), which is
why HIP-0901 routes by numeric tier instead of weakening the check. Where this path is available it
should be preferred to everything else in this section, and the confidential-compute tier should be
understood as what one uses when it is not.

The idea of pairing an enclave with Freivalds is not new: Slalom \cite{slalom} outsources linear
layers from a trusted enclave to an untrusted GPU and verifies them with exactly this check. The
composition proposed here differs in placement --- the enclave is on the operator's machine rather
than the user's, and the verifier is the settlement layer rather than the enclave --- but the
asymmetry it exploits is the same one.

\subsection{Reference-model audit}

The platform holds the advertised weights. It can therefore check a sample of served requests
directly, and the check does not involve the operator's hardware, firmware, or claims.

\paragraph{Cost of checking versus cost of serving.}
Generating $N$ tokens costs $N$ sequential forward passes, each reading the full parameter set from
memory: traffic on the order of $N \cdot P \cdot b$ bytes for $P$ parameters at $b$ bytes each, an
arithmetic intensity of about $2/b$ FLOP per byte at batch one. Checking the same completion by
teacher forcing costs one forward pass over $N$ positions: the same $\approx 2PN$ FLOPs of weight
arithmetic, but weight traffic of $P \cdot b$ rather than $N \cdot P \cdot b$, and one dependency
chain rather than $N$. Verification is not cheaper in arithmetic; it is cheaper in time and energy
by up to the roofline ratio of the hardware, which is large in the memory-bound decode regime. The
comparison holds while $N$ is small enough that attention, whose cost grows as $N^2$ in the checking
pass, does not dominate the weight term.

\paragraph{Exact replay.}
With the sampling seed committed (Section~\ref{sec:receipt}), the served output is a deterministic
function of weights, input, and seed. The audit is then an equality check. Its precondition is
arithmetic reproducibility, which does not hold by default: reduction order in a fused kernel
typically depends on the batch shape, so the same model on the same input yields different logits
under different concurrent load. Reproducible audit therefore requires batch-invariant kernel
selection on the audit path, declared in $\kappa$. This is an engineering requirement with a known
solution and a throughput cost; it is not a research problem.

\paragraph{Statistical test, where exact replay is unavailable.}
Let $p_\theta$ be the advertised model's next-token distribution at the served temperature. For a
returned completion $y_{1:N}$, define per position $t$
\[
\ell_t = \log p_\theta(y_t \mid x, y_{<t}), \qquad
H_t = -\sum_v p_\theta(v \mid \cdot) \log p_\theta(v \mid \cdot), \qquad
\sigma_t^2 = \operatorname{Var}_{v \sim p_\theta}\!\left[\log p_\theta(v \mid \cdot)\right],
\]
all computable by the auditor from the reference weights in the single forward pass above. Under the
hypothesis that the completion was served by $p_\theta$, the terms $\ell_t + H_t$ form a
martingale difference sequence with respect to the prefix filtration, and
\[
Z_N \;=\; \frac{\sum_{t=1}^{N} (\ell_t + H_t)}{\sqrt{\sum_{t=1}^{N} \sigma_t^2}}
\]
is asymptotically standard normal, subject to a Lindeberg condition that a minimum-probability floor
in the served distribution supplies. If the completion was instead served by some $q$, the per-token
drift is
\[
\delta_t \;=\; H(p_t) - H(q_t) - \mathrm{KL}(q_t \,\|\, p_t),
\]
so to first order the test's power grows with $\left|\sum_t \delta_t\right| / \sqrt{\sum_t
\sigma_t^2}$; the normalization is the null variance, and the true variance under $q$ differs, so
that expression estimates power rather than stating it.

Three limitations. The construction assumes the completion was sampled at a positive temperature:
at temperature zero the served distribution is a point mass, $H_t$ and $\sigma_t$ vanish, and the
statistic is undefined --- greedy decoding is the case for exact replay, not for this test. The test
is two-sided, because a substituted model with lower entropy can produce anomalously \emph{high}
likelihood under $p_\theta$, so a one-sided test misses an entire family of substitutions. And the
statistic has a blind set: any $q$ with $H(q_t) = H(p_t) - \mathrm{KL}(q_t\|p_t)$ at every position
is invisible to it. Additional statistics --- the rank of $y_t$ under $p_t$, or a second moment ---
shrink the blind set without emptying it. This is a detector with a false-positive rate and a power
curve, not a proof, and it should be described that way wherever it is used.

\paragraph{Indistinguishable audit traffic.}
A detector the operator can recognise has no power, because the rational response is to serve
honestly exactly when observed. Audit requests must therefore be indistinguishable from real ones:
drawn from the same distribution, issued through the same client identities and routing, and
--- critically --- opaque to the operator, which is precisely what the confidentiality of
Section~\ref{sec:flow} provides. The two mechanisms compose: confidentiality is what makes the audit
undetectable, and the audit is what detects a failure of confidentiality's own guarantees. Neither
is redundant.

Audit traffic that carries real user prompts, on the other hand, discloses those prompts to the
auditor. Synthetic requests avoid that and are the default; auditing a real request requires the
user's consent or an auditor that is itself a measured environment.

\paragraph{The deterrence condition.}
Let $\Delta$ be the cost an operator saves per substituted request, $\alpha$ the fraction of requests
audited, $\beta$ the probability that an audited substitution is detected, and $S$ the bond forfeited
on detection. Substituting has expected value $\Delta - \alpha\beta S$, so deterrence requires
\[
\alpha \beta S \;>\; \Delta .
\]
Because $\Delta$ is a marginal serving cost and $S$ is a bond, $\alpha$ can be small. The assumptions
matter more than the inequality: the operator is risk-neutral, cannot distinguish audited requests,
forfeits with certainty on detection, and has not already withdrawn the bond. The condition also
bounds expected value and not damage --- the number of substitutions before the first detection is
geometric with mean $(\alpha\beta)^{-1}$, so where harm is not linear in $\Delta$, $S$ must cover
that run and not merely one request. HIP-0095's challenge cadence and HIP-0096's bond sizing are
where these numbers are set; this document supplies the inequality they must satisfy, not the
numbers.

\subsection{Refereed delegation}

Where two independent operators serve the same request, a disagreement can be resolved by bisection
to the first divergent step, with the loser slashed \cite{crr2011}. This is the mechanism behind
optimistic fraud proofs and it gives a strong guarantee under one honest operator. It costs a second
execution and requires bit-reproducible arithmetic across independent hardware, which is a stronger
requirement than batch invariance on one machine and is not currently met by heterogeneous GPU
fleets. It is listed for completeness, not proposed.

\subsection{Succinct proofs}

A succinct argument over the whole forward pass would establish execution integrity outright, with
no hardware assumption. Current proving costs place it beyond production language-model inference by
orders of magnitude; HIP-0901 assigns it to small models and decisions, which is the right scope for
it today.

\section{The boundary}
\label{sec:boundary}

The classification below is the document's actual conclusion. \emph{C} is cryptographic --- it holds
against a computationally unbounded or bounded adversary by a stated hardness assumption, with no
trust in a manufacturer. \emph{H} is hardware trust --- it holds if named vendors, firmware, and
measured software behave as specified. \emph{E} is economic --- it holds if a rational adversary
finds cheating unprofitable. \emph{O} is open.

\begin{center}
\small
\begin{tabular}{@{}L{0.30\textwidth}L{0.30\textwidth}cL{0.26\textwidth}@{}}
\toprule
Property & Mechanism & Class & What breaks it \\
\midrule
Prompt confidential in transit & X-Wing HPKE sealing to $pk_e$ & C & ML-KEM-768 \emph{and} X25519
both broken \\
Prompt confidential during CPU execution & CPU TEE memory encryption and isolation & H & Vendor key,
firmware bug, side channel, physical attack \\
Prompt confidential during GPU execution & GPU confidential mode (tier $\geq 3$) & H & As above, plus
GPU firmware; \emph{no protection at tier 2} \\
Environment ran a specific image & Attestation with measurement allowlist & H & Vendor key or
firmware; meaningless without a reproducible build \\
Report is not a replay & Verifier nonce in report data; ephemeral keys & C & Nothing, if the nonce is
checked; \emph{everything, if it is not} \\
Weights released only to a measured environment & KMS policy on verified evidence & H & Same base as
attestation, plus KMS policy correctness \\
Model commitment names the advertised weights & Merkle root over loaded bytes & C & Second preimage
on keccak \\
\textbf{The committed weights were the operands} & Receipt signed inside the environment & \textbf{H}
& Any of (H1)--(H5); \emph{$sk_r$ extraction is undetectable} \\
The committed weights were the operands, integer path & Freivalds over the opened layer & C &
Nothing; soundness is information-theoretic \\
Substitution is unprofitable & Sampled audit, bond, slash & E & Distinguishable audits; bond smaller
than the run before detection \\
\textbf{Execution integrity, floating-point path} & --- & \textbf{O} & --- \\
Metadata privacy (timing, lengths, rates) & --- & O & --- \\
Availability and consistency & Reputation, HIP-0095 challenges & E & --- \\
\bottomrule
\end{tabular}
\end{center}

Three rows deserve restating outside the table.

The row that most deployments would label ``verified inference'' is class H, not C. It is the claim
that the committed weights were actually multiplied, and it rests on the full chain
(H1)--(H5). Where the numeric tier permits, the row below it replaces H with C, and that is the only
place in this system where the replacement is available.

Execution integrity on the floating-point path is an open problem rather than an expensive or
partially solved one. Confidential computing substitutes hardware trust for it, sampled audit
substitutes economics for it, and neither is a proof. Any document, this one and our own HIPs
included, that describes the confidential tier as ``verifiable inference'' without that
qualification is overstating it.

Tier 2 gives no confidentiality for the inference itself when a discrete GPU is involved. It is
listed as a tier because it is a real step up for CPU-side work, and it should never be sold as
protecting a prompt during GPU serving.

\section{Open problems}

\paragraph{Execution integrity in floating point.}
The gap between the integer path, where Freivalds gives an information-theoretic guarantee, and the
floating-point path, where nothing does, is the central unsolved problem. Two directions are worth
work: extending the exact-arithmetic path to cover more of production serving, which is a
quantization and kernel question rather than a cryptographic one; and sound verification under
controlled non-determinism, where the difficulty is that any tolerance wide enough to admit honest
reductions is wide enough to hide a cheat, and the published proposals in this area are detectors
with error rates rather than proofs.

\paragraph{Detecting key extraction.}
$sk_r$ compromise is currently undetectable by construction. The distinct-measurement quorum raises the cost
and the audit detects the resulting behaviour, but no mechanism detects the extraction itself.
Whether periodic re-attestation with hardware-bound freshness can be strengthened into detection ---
rather than merely into a bounded window --- is open.

\paragraph{Physical attacks by the party being attested.}
The standard threat model for confidential computing does not include an adversary who owns the
machine. This deployment does. Whether any composition of remote attestation, tamper evidence, and
economic bonding meaningfully raises the cost of a physical attack, or whether the honest answer is
that certain models simply do not go on third-party hardware, is unresolved and should be decided by
policy rather than assumed away.

\paragraph{Metadata.}
Nothing here protects timing, lengths, or rates. Constant-rate token emission and length padding are
implementable at a latency cost that has not been measured.

\paragraph{Post-quantum attestation.}
Hardware attestation signatures are classical elliptic-curve or RSA today. \code{TEEEvidence}
already records this as the sole classical assumption in an otherwise post-quantum tier, and notes
that it is governance-removable and that a post-quantum device-identity path would remove it. Until
such roots exist in silicon, the attestation leg of an otherwise post-quantum settlement path is not
post-quantum, and should not be described as such.

\section{What is missing in the tree}

Ordered by what blocks what, and stated as work rather than as a plan.

There must be one attester and one verifier, and there are currently three places that could hold
them. The mocks in \code{net/hanzo-pqc} and \code{rust-sdk/crates/hanzo-kbs} return success for any
input; \code{hanzo-attest} (Section~\ref{sec:hanzoattest}) declares the two hardware roots and
refuses them by name, which is where the real implementation belongs, because refusing by name is
already the correct behaviour for an unimplemented root and the mocks are not. That implementation
performs the checks of Section~\ref{sec:evidence} --- certificate chain to vendor root, measurement
allowlist, policy bits, TCB floor, report-data binding, freshness nonce --- and is consumed by the
key broker and by the chain-side voucher alike. The two overlapping tier enumerations collapse to one
at the same time.

The verified payload must carry the computation binding. Section~\ref{sec:precompile} observed that
the 48-byte Lux header has no report-data field, so the third check \code{TEEEvidence} documents
cannot be performed. Either the verifier parses genuine SEV-SNP and TDX reports and reads their
64-byte field, or the header gains one; the first is preferable because it also supplies the
measurement, policy, and TCB fields the checklist needs, none of which the header carries either.

The vendor root pool holds one root, Intel's. AMD's ARK and NVIDIA's device-identity root are needed
before an SEV-SNP host or a confidential GPU can attest at all.

Measurements must be reproducible and published. Without the build-to-measurement mapping and its
log, the governance-approved measurement set is a list of values nobody outside can interpret, and
requiring a quorum over them adds no assurance.

The receipt of Section~\ref{sec:receipt} is not emitted anywhere. The engine's transcript machinery
(\code{poi\_transcript.rs}, \code{poi\_forward.rs}) is the natural place, and HIP-0901 already
records that emitting a transcript from a live forward pass is its remaining engineering.

The sampled audit does not exist: no canary traffic, no reference-model check, no batch-invariant
audit path. It is also the only mechanism in this document that detects a compromised environment
rather than assuming one cannot occur, which argues for building it early rather than last.

Finally, the confidential-VM backend does not exist in \code{hanzoai/vm}. The microVM layer boots
k3s quickly on a normal hypervisor; a confidential guest is a different launch path with a different
firmware and a measurement to reproduce, and none of the flow above can be exercised until it lands
on a host that has the hardware --- which, per Section~\ref{sec:observed}, the bench currently does
not.

\begin{thebibliography}{99}

\bibitem{freivalds1977}
R. Freivalds.
\newblock Probabilistic machines can use less running time.
\newblock \emph{IFIP Congress}, 1977.

\bibitem{slalom}
F. Tram\`er and D. Boneh.
\newblock Slalom: Fast, verifiable and private execution of neural networks in trusted hardware.
\newblock \emph{International Conference on Learning Representations}, 2019.

\bibitem{crr2011}
R. Canetti, B. Riva, and G. N. Rothblum.
\newblock Practical delegation of computation using multiple servers.
\newblock \emph{ACM Conference on Computer and Communications Security}, 2011.

\bibitem{buhren2021}
R. Buhren, H.-N. Jacob, T. Krachenfels, and J.-P. Seifert.
\newblock One glitch to rule them all: Fault injection attacks against AMD's Secure Encrypted
Virtualization.
\newblock \emph{ACM Conference on Computer and Communications Security}, 2021.

\bibitem{cipherleaks}
M. Li, Y. Zhang, H. Wang, K. Li, and Y. Cheng.
\newblock CipherLeaks: Breaking constant-time cryptography on AMD SEV via the ciphertext side
channel.
\newblock \emph{USENIX Security Symposium}, 2021.

\bibitem{sevstep}
L. Wilke, J. Wichelmann, A. Rabich, and T. Eisenbarth.
\newblock SEV-Step: A single-stepping framework for AMD-SEV.
\newblock 2024.

\bibitem{heckler}
B. Schl\"uter, S. Sridhara, A. Bertschi, and S. Shinde.
\newblock Heckler: Breaking confidential VMs with malicious interrupts.
\newblock 2024.

\bibitem{wesee}
B. Schl\"uter, S. Sridhara, M. Kuhne, A. Bertschi, and S. Shinde.
\newblock WeSee: Using malicious \#VC interrupts to break AMD SEV-SNP.
\newblock 2024.

\bibitem{cachewarp}
R. Zhang, L. Gerlach, D. Weber, L. Hetterich, Y. Lu, A. Kogler, and M. Schwarz.
\newblock CacheWarp: Software-based fault injection using selective state reset.
\newblock \emph{USENIX Security Symposium}, 2024.

\bibitem{badram}
J. De Meulemeester, L. Wilke, D. Oswald, T. Eisenbarth, I. Verbauwhede, and J. Van Bulck.
\newblock BadRAM: Practical memory aliasing attacks on trusted execution environments.
\newblock 2025.

\bibitem{snpabi}
Advanced Micro Devices.
\newblock \emph{SEV Secure Nested Paging Firmware ABI Specification}.

\bibitem{tdx}
Intel Corporation.
\newblock \emph{Intel Trust Domain Extensions (TDX) Module Architecture Specification} and
\emph{Data Center Attestation Primitives}.

\bibitem{nvcc}
NVIDIA Corporation.
\newblock \emph{Confidential Computing Deployment Guide} and \emph{Attestation Services
Documentation}.

\bibitem{hpke}
R. Barnes, K. Bhargavan, B. Lipp, and C. Wood.
\newblock Hybrid Public Key Encryption.
\newblock RFC 9180, 2022.

\bibitem{rfc6962}
B. Laurie, A. Langley, and E. Kasper.
\newblock Certificate Transparency.
\newblock RFC 6962, 2013.

\bibitem{hip0901}
Hanzo AI.
\newblock HIP-0901: Proof of AI --- native execution proofs, canonical contract, and operator-LLM
governance.

\bibitem{hip0095}
Hanzo AI.
\newblock HIP-0095: QoS challenge system.

\bibitem{hip0096}
Hanzo AI.
\newblock HIP-0096: AI compute contribution rewards.

\bibitem{hip1325}
Hanzo AI.
\newblock HIP-1325: Node --- a machine an org owns.

\bibitem{hanzodattest}
Hanzo AI Research.
\newblock Attested compute in \code{hanzod}: boot measurement, host-asserted evidence, and the gap to
hardware attestation.
\newblock Hanzo Papers.

\end{thebibliography}

\end{document}
